%=============================================================================!
% 10.6  THE LOOP                                                              !
%=============================================================================!
\section{The loop}
\label{sec:md-loop}

The pieces now fit together: a model that returns the energy and forces
of atoms in a cell, the velocity Verlet step of \cref{sec:in-verlet}, and
a record of what happens. This section writes the loop, checks it against
an independent library, and runs it on argon and on a model salt.

\subsection*{A model and a loop}

A model is anything that takes the positions and returns the potential
energy and the forces. For a pair potential in a cell, \texttt{PairModel}
of \texttt{mdlab.md} gets the pairs within $r_{\mathrm c}$ from its Verlet
list and adds up their forces:

\begin{code}
    def __call__(self, positions: ArrayLike) -> tuple[float, NDArray]:
        r = np.asarray(positions, dtype=float)
        i, j, d, dist = self.neighbours.pairs(r)
        phi, dphi = self.pair(dist)
        push = (dphi / dist)[:, None] * d  # the force on i from j
        forces = np.zeros_like(r)
        np.add.at(forces, i, push)
        np.add.at(forces, j, -push)
        return float(np.sum(phi)), forces
\end{code}

\noindent With $\vb{d} = \vb{r}_j - \vb{r}_i$, the force on atom $i$ from
the pair is $\varphi'(r)\,\vb{d}/r$ (\cref{sec:pe-forces}), and the force on
$j$ its opposite; \texttt{np.add.at} adds each pair's push into the rows
of its two atoms, correctly when an atom appears in many pairs.
\texttt{EwaldModel} returns, in the same form, the energy and forces of
charges by the Ewald sum of \cref{sec:md-ewald}, and \texttt{sum\_of}
adds the energies and forces of several models.

The loop of \texttt{run} is velocity Verlet over any such model:

\begin{code}
    u, f = model(r)
    for step in range(n_steps + 1):
        if step > 0:
            v += 0.5 * dt * force_to_accel * f / m
            r += dt * v
            u, f = model(r)
            v += 0.5 * dt * force_to_accel * f / m
        if step % every == 0:
            kinetic = 0.5 * to_energy * float(np.sum(m * v * v))
\end{code}

\noindent Each step is a half kick, a drift and a half kick, with one call
of the model, whose forces serve the next step's first half kick, as in
\cref{fig:in-kdk}. Here \texttt{m} holds the masses as a column, one row
per atom, and \texttt{to\_energy} is \texttt{MV2\_TO\_EV}, or 1 in
reduced units. Every \texttt{every} steps the time, positions, velocities
and both energies are recorded. If a file is named, the frames are written
to it in the extended XYZ format, which ASE reads: for each frame, the
number of atoms; a line with the lattice vectors, the names of the columns
and the potential and kinetic energies; and a line per atom with its
element, position and velocity. ASE takes a value named \texttt{energy}
as the potential energy, as the tests of \texttt{mdlab} check, so the
potential energy is written under that name.

\subsection*{Checked against ASE}

ASE, the Atomic Simulation Environment\cite{Larsen2017}, has its own
Lennard-Jones forces and its own velocity Verlet. For 256 argon atoms in
a face-centred cubic cell, with the energy shifted at the cutoff in both,
the two give the same starting energy to every printed figure. After 200
steps of \SI{5}{\femto\second} the positions at first differed by
$\num{1.6e-8}$\,\AA. The cause lay in the units: ASE converts
femtoseconds into its own unit of time with the 2014 values of the
physical constants, and \texttt{mdlab} with the 2022 values that SciPy
ships, and the two femtoseconds differ by 4.6 parts in $10^9$. An argon
atom carrying \SI{0.02}{\electronvolt} moves at about
\SI{0.0031}{\angstrom\per\femto\second}, and over \SI{1000}{\femto\second}
a clock off by that fraction misplaces it by about $0.0031\times1000
\times\num{4.6e-9} = \num{1.4e-8}$\,\AA, the size found
(\cref{ex:md-femtosecond}). Given the same femtosecond, the positions after
200 steps agree to $\num{1.8e-14}$\,\AA\ and the velocities to
$\num{1.1e-16}$\,\AA/fs.

\begin{practice}
A comparison of two codes needs the same units, constants and starting
state, and only a few hundred steps: the motion of many atoms is chaotic,
two paths that start a little apart separating faster and faster, and any
difference, even of rounding, grows over a long run until the paths share
nothing but their statistics\cite{FrenkelSmit2023}. A difference that
grows steadily from the first step, as this one did, from
$\num{6e-13}$\,\AA\ after 2 steps to $\num{4.8e-9}$\,\AA\ after 50 and
$\num{1.6e-8}$\,\AA\ after 200, points to something systematic before
the chaos takes over.
\end{practice}

\subsection*{Two runs}

\Cref{fig:md-run}(a) follows 500 argon atoms (\cref{sec:md-start}) for
\SI{20}{\pico\second} with $\dt = \SI{10}{\femto\second}$, the
Lennard-Jones energy switched off between $2\sigma$ and $2.5\sigma$ and a
skin of \SI{1}{\angstrom}. The total energy stays within
$\num{3.7e-5}$\,eV per atom, and over the last \SI{10}{\pico\second} its
standard deviation is 0.3\% of that of the kinetic energy, the test of
\cref{sec:in-stability}. About \SI{10}{\milli\electronvolt} per atom
passes from kinetic to potential energy in the first half picosecond,
after which the kinetic energy swings by \SI{1.6}{\milli\electronvolt};
the total momentum at the end of the run is that of
\cref{sec:md-periodic}, the rounding of the arithmetic.

\begin{figure}[htb]
\centering
\includegraphics{ch10_code/run}
\caption{Two periodic runs. (a) 500 argon atoms started on their
face-centred cubic lattice with \SI{0.02}{\electronvolt} of kinetic energy
per atom, $\dt = \SI{10}{\femto\second}$: the kinetic energy, the potential
energy and the total, each per atom and the last two less the starting
potential energy, over the first \SI{5}{\pico\second} of
\SI{20}{\pico\second}. (b) The model salt of 64 ions with Ewald
electrostatics: the change of the total energy over $2\tau$ with $\dt =
0.002\tau$ and $0.004\tau$.}
\label{fig:md-run}
\end{figure}

\Cref{fig:md-run}(b) runs a model salt, chosen to exercise the Ewald sum
in a loop. Its 64 ions, of charge $\pm1$ and mass 1 in reduced units, sit
on a rock-salt lattice, repel each other at contact through the
Lennard-Jones energy cut off at its minimum, $2^{1/6}\sigma$, and shifted
there to zero, and attract and repel through the Coulomb energy with
$k_{\mathrm e} = 10\,\varepsilon\sigma$, a value chosen for this model,
summed by \texttt{EwaldModel}. Its energy is least with neighbouring ions
$1.0908\sigma$ apart, at $-15.810\varepsilon$ per ion pair, found as for
argon. Given $0.5\varepsilon$ of kinetic energy per ion and followed for
$2\tau$, the total energy of the 64 ions spreads by $\num{6.08e-3}\,\varepsilon$ with $\dt =
0.002\tau$ and by $\num{2.46e-2}\,\varepsilon$ with $0.004\tau$, a ratio of
4.04 for a doubled step: the $\dt^2$ of \cref{sec:in-stability}, the
test that the Ewald forces are the slopes of the Ewald energy and are
smooth.

\begin{keyidea}
A model returns the energy and the forces; the loop is velocity Verlet
over it with one call per step. The code agrees with ASE to the rounding
of the arithmetic once both use the same constants, keeps the energy of a
periodic crystal with a fluctuation 0.3\% of that of its kinetic energy
and its momentum to rounding, and its Ewald forces pass the $\dt^2$ test.
\end{keyidea}

\notebook{10}{plot the argon and model-salt runs, check the loop against ASE, and read a trajectory file back}

\begin{selfcheck}
With $\dt = \SI{10}{\femto\second}$, how many steps are 20\,ps, and how
many force calls did the argon run make, counting the first?
\end{selfcheck}
